\section{O Protótipo Industrial Entregue}
\label{sec:resultado_prototipo}

Para além do desempenho das arquiteturas neurais, o desenvolvimento do \textit{hardware} de captura constitui um resultado central desta dissertação. O protótipo final, ilustrado nas fotografias da Figura \ref{fig:caixa_real}, provou ser resiliente ao ambiente de colheita, mantendo a estabilidade fotométrica necessária para que o modelo ConvNeXt operasse com a precisão reportada. 

A autonomia do equipamento permitiu a digitalização de mais de 1.200 amostras reais com zero índice de falha mecânica, demonstrando que a integração entre a câmara de luz controlada e a unidade de processamento de borda é tecnicamente viável para a escala industrial do cerrado brasileiro.

\begin{figure}[H]
    \centering
    \caption{Fotografias reais do protótipo industrial: (a) visão detalhada do hardware e (b) demonstração de uso no fluxo de trabalho industrial.}
    \label{fig:caixa_real}
    % Primeira imagem (Caixa em detalhe) - Recortada embaixo para remover espaço vazio
    \begin{subfigure}[b]{0.48\textwidth}
        \centering
        \includegraphics[trim={20 140 20 30}, clip, width=0.93\textwidth]{img_caixa}
        \caption{Detalhe do hardware e compartimento interno.}
        \label{fig:caixa_detalhe}
    \end{subfigure}
    \hfill
    % Segunda imagem (Operação) - Recortada apenas embaixo para manter o enquadramento superior
    \begin{subfigure}[b]{0.48\textwidth}
        \centering
        \includegraphics[trim={0 200 0 50}, clip, width=\textwidth]{img_caixa_operacao}
        \caption{Protótipo em operação real na algodoeira.}
        \label{fig:caixa_operacao}
    \end{subfigure}
    
    \legend{Fonte: Do autor (2026).}
\end{figure}
